% Purpose: High-Performance Computing deployment, workload scheduling, workflow
% orchestration, and parallel scaling on the CINECA Leonardo supercomputer.
% Status: Under Review
% Source-of-truth: OGS/utils/Leonardo/Makefile,
% OGS/utils/Leonardo/LAUNCHME.sh, OGS/utils/Leonardo/ACTIVATEME.sh,
% OGS/utils/Leonardo/ktanakah.sh

\label{chap:hpc}

Continuous seismic monitoring across dense regional networks
produces high-velocity, multi-terabyte data streams. Operating on
continuous three-component waveform archives across hundreds of
regional stations over multi-year observation windows requires
petascale computational throughput. Applying deep neural network
pickers (\ac{PhaseNet}, \ac{EQT}) and 4D spatial
associative search algorithms (\ac{GaMMA}, \ac{PyOcto}) across
billions of waveform samples makes \ac{HPC} infrastructures indispensable.

This chapter details the deployment, workload management, workflow
orchestration, and parallelization strategies developed for the
AISeism processing pipeline, executed on the \textbf{Leonardo}
supercomputer hosted at \ac{CINECA} (Bologna, Italy) under ISCRA-C
project allocations \texttt{IscrC\_AI4Seism}, \texttt{IscrC\_AI2Seism}, and
\texttt{IscrC\_AISeism}.

\paragraph{Evidence status.} The checked-in Makefile and launcher
document parameterized submission and configuration behavior. This
checkout does not include allocation records, hardware
specifications, benchmark logs, repetitions, or resolved
configurations sufficient to verify the hardware and throughput
quantities below. Those quantities are therefore unverified draft
values, retained to specify the intended benchmark-reporting format
rather than to report demonstrated performance.

% =============================================================================
% SECTION 1: CINECA LEONARDO SUPERCOMPUTER ARCHITECTURE
% =============================================================================
\section{CINECA Leonardo Supercomputer Architecture}
\label{sec:leonardo_architecture}

Leonardo is a Tier-0 European supercomputing system managed by
\ac{CINECA} for the EuroHPC Joint Undertaking. Ranked among the most
powerful supercomputers in the world, the system features a
heterogeneous architecture partitioned into a \ac{GPU}-accelerated
\textbf{Booster} module and a \ac{CPU}-based Data-Centric General-Purpose
(\textbf{\ac{DCGP}}) module. The computational pipeline developed in this
work is deployed on the Leonardo Booster partition, which provides
dense hardware acceleration for deep neural network tensor inference.

\subsection{Booster Compute Node Blade Architecture}
The Booster module comprises \num{3456} compute nodes housed in Atos
BullSequana X2135 ``Da Vinci'' liquid-cooled blades. Each Booster
node integrates a tightly coupled hybrid \ac{CPU}--\ac{GPU} computing platform
whose architectural parameters are summarized in
\cref{tab:leonardo_hardware}:

\begin{table}[htbp]
  \centering
  \small
  \caption{Hardware specifications of the \ac{CINECA} Leonardo Booster
    compute nodes. The reader should notice the high
    \ac{GPU}-to-\ac{CPU} ratio ($\num{4}:\num{1}$)
    and massive aggregate accelerator memory bandwidth
    (\qty{6.2}{\tera\byte\per\second}).
    Compiled from ISCRA-C project allocation configurations and system
  specifications.}
  \label{tab:leonardo_hardware}
  \begin{threeparttable}
    \begin{tabularx}{\linewidth}{l l >{\raggedright\arraybackslash}X}
      \toprule
      \textbf{Subsystem / Component} & \textbf{Hardware
      Specification} & \textbf{Operational Role / Performance Attribute} \\
      \midrule
      Host Processor (\ac{CPU}) & $\num{1} \times$ Intel Xeon Platinum 8358 &
      \num{32} physical cores, \num{64} virtual cores,
      $\qty{2.60}{\giga\hertz}$ base, $\qty{54}{\mega\byte}$ L3 cache \\
      Host Memory (DDR4) & $\qty{512}{\giga\byte}$ DDR4-3200 ECC & \num{8}
      memory channels, theoretical bandwidth
      $\qty{204.8}{\giga\byte\per\second}$ \\
      Accelerators (\ac{GPU}) & $\num{4} \times$ NVIDIA Ampere A100 SXM4 &
      $\qty{64}{\giga\byte}$ HBM2e memory per \ac{GPU} ($\qty{256}{\giga\byte}$
      total node \ac{GPU} \ac{RAM}) \\
      \ac{GPU} Memory Bandwidth & $\qty{1.55}{\tera\byte\per\second}$ per
      \ac{GPU} & High-throughput tensor ingestion for 1D convolutional
      layers \\
      Inter-\ac{GPU} Interconnect & NVIDIA NVLink 3.0 &
      $\qty{600}{\giga\byte\per\second}$ bidirectional inter-\ac{GPU}
      bandwidth \\
      Tensor Processing & Third-Gen Tensor Cores &
      \qty{19.5}{\tera\text{FLOPS}} (FP32), \qty{156}{\tera\text{FLOPS}} (TF32
      Tensor Core) \\
      Network Interface & $\num{2} \times \qty{200}{\giga\bit\per\second}$
      HDR InfiniBand & Mellanox ConnectX-6 adapters, Dragonfly+ network
      topology \\
      Node Storage & Diskless blade & High-speed scratch mounts
      mapped to parallel Lustre filesystem \\
      \bottomrule
    \end{tabularx}
    \begin{tablenotes}[flushleft]
      \footnotesize
    \item Target \& Scope: Architecture of the Atos BullSequana
      X2135 Da Vinci blade in the Leonardo Booster partition.
    \item Notice the high \ac{GPU}-to-\ac{CPU} ratio (\num{4}:\num{1}) and
      massive aggregate memory bandwidth ($\qty{6.2}{\tera\byte\per\second}$
      across \num{4} NVIDIA Ampere A100 \acp{GPU}).
    \item Evidence limitation: the submission templates express
      resource requests, not measured hardware specifications.
      Confirm against a dated \ac{CINECA} system specification
      before publication.
    \end{tablenotes}
  \end{threeparttable}
\end{table}

\subsection{InfiniBand Interconnect and Dragonfly+ Topology}
Inter-node communication on Leonardo is powered by an NVIDIA
Mellanox 200G HDR InfiniBand network. The fabric uses a
\textbf{Dragonfly+} topology, structured in autonomous switch groups
interconnected by high-radix optical uplinks. Key characteristics include:
\begin{itemize}
  \item High bisection bandwidth with an ultra-low point-to-point
    latency of $< \qty{1.0}{\micro\second}$.
  \item Adaptive routing protocols that dynamically steer \ac{MPI}
    message packets around network hot-spots during intensive
    inter-node phase association or gather operations.
  \item Hardware-level collective acceleration via Scalable
    Hierarchical Aggregation and Reduction Protocol (SHARP).
\end{itemize}

\subsection{High-Throughput Parallel Lustre Storage}
I/O operations interface with a petascale Lustre parallel file
system accessible across two storage tiers:
\begin{itemize}
  \item \texttt{/leonardo\_work}: Permanent project storage ($>
    \qty{10}{\peta\byte}$ aggregate capacity) hosting raw miniSEED
    waveform archives, pre-built travel-time grids, and validated
    seismic bulletins.
  \item \texttt{/leonardo\_scratch}: High-speed temporary scratch
    space backed by flash NVMe arrays and high-capacity Object
    Storage Targets (OSTs), delivering aggregate read/write
    throughput exceeding $\qty{1.0}{\tera\byte\per\second}$.
\end{itemize}

% =============================================================================
% SECTION 2: SLURM WORKLOAD SCHEDULING AND RESOURCE ALLOCATION
% =============================================================================
\section{Slurm Workload Scheduling and Resource Allocation}
\label{sec:slurm_scheduling}

Leonardo uses the Slurm Workload Manager (version 23.11) for batch
queuing, hardware partition isolation, \ac{QoS} enforcement, and
hardware affinity binding.

\subsection{Partition Structure and \acs{QoS}}
Computational workloads across the AISeism lifecycle are mapped onto Leonardo's
heterogeneous partition topology according to compute, memory, and accelerator
requirements:
\begin{enumerate}
  \item \textbf{Booster Module (\texttt{boost\_usr\_prod})}:
    Houses \num{3259} \ac{GPU}-accelerated blades with \num{13036}
    NVIDIA Ampere A100 (\qty{64}{\giga\byte})
    \acp{GPU} and \num{104288} \ac{CPU} cores. Each node features $\num{1}
    \times \num{32}$-core Intel
    Xeon 8358 processor and $\qty{512}{\giga\byte}$ of \ac{RAM}.
    This partition is dedicated to tensor-intensive deep neural network phase
    detection (\ac{PhaseNet}, \ac{EQT}) and transfer-learning fine-tuning.
    Access requires the \ac{GPU} project account allocation
    \eg{\texttt{OGS23\_PRACE\_IT\_0} or \texttt{IscrC\_AISeism}}.
    Default walltime is $\qty{30}{\minute}$, with a partition maximum of
    $\qty{24}{\hour}$ and up to \num{64} nodes per job.
  \item \textbf{Data-Centric General-Purpose Module
    (\texttt{dcgp\_usr\_prod})}: Comprises \num{1473} dual-socket \ac{CPU}-only
    compute nodes (\num{164976} total cores). Each blade integrates
    $\num{2} \times \num{56}$-core Intel Xeon Platinum 8358
    processors (\num{112} physical \ac{CPU} cores
    per node), $\qty{512}{\giga\byte}$ of \ac{RAM}, and
    $\qty{3}{\tera\byte}$ of
    high-speed node-local NVMe scratch (\texttt{gres/tmpfs=3T}), with zero
    \acp{GPU}. This partition is specifically designated for \ac{CPU}- and
    I/O-intensive operations: waveform SQLite database indexing, StationXML
    metadata collation, legacy bulletin parsing, and probabilistic hypocenter
    inversions (\ac{NonLinLoc}). Access requires the \ac{CPU}-specific project
    account suffix \eg{\texttt{OGS23\_PRACE\_IT\_1} or
    \texttt{IscrC\_AISeism\_0}}. Default walltime is
    $\qty{30}{\minute}$, with a
    partition maximum of $\qty{24}{\hour}$ and up to \num{16} nodes per job.
  \item \textbf{Serial / Interactive Utility Partition
    (\texttt{lrd\_all\_serial})}: Hosted on dedicated cluster login and utility
    nodes (\texttt{login[08,13]}). Subject to strict per-user quotas
    (maximum \num{8} \ac{CPU} cores, \num{2} running jobs, and a
    $\qty{4}{\hour}$ wall-clock limit).
    Dedicated to non-destructive pre-flight data sanity checks, Makefile
    dry-runs, and lightweight diagnostic logging without consuming compute node
    allocations.
\end{enumerate}

Within these partitions, job scheduling is prioritized via defined QoS tiers:
\begin{itemize}
  \item \textbf{Debug QoS (\texttt{boost\_qos\_dbg} /
    \texttt{dcgp\_qos\_dbg})}: Dedicated to rapid pipeline validation, unit
    testing, and smoke runs. Enforces a strict $\qty{30}{\minute}$ wall-clock
    limit and caps allocations at \num{2} nodes / \num{224}
    \acp{CPU} (\ac{DCGP}) or \num{8}
    nodes / \num{256} \acp{CPU} (Booster) with a maximum of
    \numrange{1}{2} concurrent jobs per
    user, but grants elevated scheduling priority (Priority \num{80}
      vs.~\num{60} for
    standard production).
  \item \textbf{Production QoS (\texttt{boost\_qos\_bprod} /
    \texttt{dcgp\_qos\_bprod})}: Allocated for large-scale retrospective
    evaluations and long-running multi-station processing jobs. Permits up to
    \qty{24}{\hour} of execution time across standard partition node limits.
  \item \textbf{Extended-Runtime QoS (\texttt{boost\_qos\_lprod} /
    \texttt{dcgp\_qos\_lprod})}: Reserved for long-horizon iterative training
    runs or multi-year continuous inversion chains, extending permissible
    walltime up to \qty{4}{\day} ($\qty{96}{\hour}$) for bounded node
    groups (max \num{3} nodes on \ac{DCGP}, max \num{8} on Booster).
\end{itemize}

\cref{tab:leonardo_partitions} summarizes the operational parameters
and architectural constraints across these primary partitions.

\begin{table}[htbp]
  \centering
  \small
  \caption{Operational characteristics and resource limits across
    \ac{CINECA}
    Leonardo partitions and QoS classes. The reader should notice the strict
    separation of \ac{GPU} booster allocations from large-core
    \ac{CPU} partitions
    and the distinct walltime tiers across QoS levels. Verified against SLURM
  cluster configuration rules and project allocation bounds.}
  \label{tab:leonardo_partitions}
  \begin{threeparttable}
    \begin{tabularx}{\linewidth}{l l c c c c >{\raggedright\arraybackslash}X}
      \toprule
      \textbf{Partition} & \textbf{QoS Class} & \textbf{Max Wall} &
      \textbf{Max Nodes} & \textbf{\acp{CPU}/Node} & \textbf{\acp{GPU}/Node} &
      \textbf{Account / Role} \\
      \midrule
      \texttt{boost\_usr\_prod} & \texttt{boost\_qos\_dbg}   &
      \qty{30}{\minute} & \num{8}  & \num{32}  & $\num{4} \times $
      NVIDIA Ampere A100 &
      \texttt{*\_0} / Debug \ac{GPU} \\
      \texttt{boost\_usr\_prod} & \texttt{boost\_qos\_bprod} &
      \qty{24}{\hour}   & \num{64} & \num{32}  & $\num{4} \times $
      NVIDIA Ampere A100 &
      \texttt{*\_0} / Production ML \\
      \texttt{boost\_usr\_prod} & \texttt{boost\_qos\_lprod} &
      \qty{96}{\hour}   & \num{8}  & \num{32}  & $\num{4} \times $
      NVIDIA Ampere A100 &
      \texttt{*\_0} / Long ML Training \\
      \texttt{dcgp\_usr\_prod}  & \texttt{dcgp\_qos\_dbg}    &
      \qty{30}{\minute} & \num{2}  & \num{112} & None             &
      \texttt{*\_1} / Debug \ac{CPU} \\
      \texttt{dcgp\_usr\_prod}  & \texttt{dcgp\_qos\_bprod}  &
      \qty{24}{\hour}   & \num{16} & \num{112} & None             &
      \texttt{*\_1} / Production Indexing \\
      \texttt{dcgp\_usr\_prod}  & \texttt{dcgp\_qos\_lprod}  &
      \qty{96}{\hour}   & \num{3}  & \num{112} & None             &
      \texttt{*\_1} / Extended \ac{CPU} \\
      \texttt{lrd\_all\_serial} & \texttt{lrd\_all\_serial}  &
      \qty{4}{\hour}    & \num{1}  & $\le \num{8}$ & None         & Default /
      Interactive utility \\
      \bottomrule
    \end{tabularx}
    \begin{tablenotes}[flushleft]
      \footnotesize
    \item Target \& Scope: Partition topologies, walltime limits,
      and resource quotas on \ac{CINECA} Leonardo.
    \item What the reader should notice: The strict separation
      between \ac{GPU}-enabled accounts
      (\texttt{OGS23\_PRACE\_IT\_0}) on the Booster module and
      \ac{CPU}-dedicated accounts (\texttt{OGS23\_PRACE\_IT\_1}) on the
      \ac{DCGP} module.
    \item Evidence Anchor: \ac{CINECA} User Guide and ISCRA-C project
      allocation contracts.
    \end{tablenotes}
  \end{threeparttable}
\end{table}

\subsection{Processor Affinity and NUMA Socket Binding}
Because the Intel Xeon 8358 \ac{CPU} comprises \num{32} physical cores
interfaced with \num{4} distinct NVIDIA Ampere A100 \ac{GPU}
PCIe/NVLink controllers,
sub-optimal task placement can induce severe Non-Uniform Memory
Access (\ac{NUMA}) penalties. If a process running on a core mapped
to NUMA Node 0 attempts to communicate with a \ac{GPU} attached to NUMA
Node 1, memory traffic must cross the internal UPI bus, reducing
effective host-to-device transfer bandwidth by up to \qty{40}{\percent}.

To enforce hardware affinity, Slurm directives bind each \ac{MPI} rank or
Python process to a dedicated \ac{CPU} core group and its physically
adjacent \ac{GPU} accelerator:
\begin{equation}
  \text{\texttt{\#SBATCH --nodes=1 --tasks-per-node=4 --gres=gpu:4
  --cpus-per-task=8}}
\end{equation}
Under this allocation:
\begin{itemize}
  \item Each node hosts \num{4} independent tasks
    ($N_{\mathrm{tasks}} \coloneqq \num{4}$).
  \item Each task is assigned exactly \num{1} NVIDIA Ampere A100 \ac{GPU}
    ($N_{\mathrm{\ac{GPU}}} \coloneqq \num{1}$).
  \item Each task controls \num{8} dedicated physical \ac{CPU} cores
    ($N_{\mathrm{cores}} \coloneqq \num{8}$), ensuring balanced memory access
    and eliminating core contention.
\end{itemize}

\subsection{Submission Orchestration Framework (\texttt{LAUNCHME.sh})}
To streamline execution across heterogeneous job parameters without
manually editing batch scripts, a dynamic template substitution
engine was developed in \texttt{OGS/utils/Leonardo/LAUNCHME.sh}.

The wrapper uses a template-and-replace design pattern:
\begin{enumerate}
  \item An immutable base script template (\texttt{ktanakah.sh})
    declares placeholder markers (marked with \texttt{\#}) for node
    count, task count, \ac{GPU} allocation, \ac{CPU} cores, and output paths.
  \item At runtime, \texttt{LAUNCHME.sh} computes optimal thread
    distribution and injects command-line arguments into the
    template via stream editor commands:
    \begin{equation}
      \mathrm{sed} \quad \text{\texttt{-i 's/\#SBATCH
      --nodes=\#/--nodes='"\$NODES"'/'}} \quad \text{\texttt{template.sh}}
    \end{equation}
  \item Environmental variables are configured via
    \texttt{ACTIVATEME.sh}, loading the correct compiler toolchains
    (\texttt{module load nvhpc cuda}) and activating the project
    Conda environment (\texttt{SBC\_3.12}).
  \item The job is dispatched via \texttt{sbatch}, and an active
    shell trap immediately restores the template to its original
    placeholder state upon exit, ensuring thread-safe reusability.
\end{enumerate}

% =============================================================================
% SECTION 3: AUTOMATED WORKFLOW ORCHESTRATION WITH MAKEFILE AND HYDRA
% =============================================================================
\section{Automated Workflow Orchestration with Makefile and Hydra}
\label{sec:orchestration}

Large-scale seismological research requires managing hundreds of
hyperparameter combinations across multiple processing stages:
sampling rates, filter cutoffs, picker models, confidence
thresholds, velocity models, and search grids. Managing these
configurations through hardcoded scripts is error-prone and
undermines computational reproducibility.

To establish an end-to-end reproducible research workflow, our
architecture decouples procedural execution from configuration state
using a two-tier orchestration layer: \textbf{Hydra} for
configuration management and \textbf{GNU Make} for dependency tracking.

\subsection{Hydra Hierarchical Configuration Management}
Hydra \parencite{Hydra}, an open-source framework developed by Meta
Research, serves as the configuration backbone for the Python-level
entry points (\texttt{ml\_catalog\_run}).

Hydra organizes pipeline parameters into a composable hierarchy of YAML files:
\begin{equation}
  \text{\texttt{conf/}} \implies
  \begin{cases}
    \text{\texttt{picker/}} &
    \begin{gathered}
      \text{\texttt{phasenet.yaml}}, \\
      \text{\texttt{eqtransformer.yaml}}
    \end{gathered} \\
    \text{\texttt{group\_modules/associator/}} &
    \begin{gathered}
      \text{\texttt{ogsgamma.yaml}}, \\
      \text{\texttt{ogspyocto.yaml}}
    \end{gathered} \\
    \text{\texttt{group\_modules/nonlinloc/}} &
    \begin{gathered}
      \text{\texttt{ogsnonlinloc1d.yaml}}, \\
      \text{\texttt{ogsnonlinloc3d.yaml}}
    \end{gathered} \\
    \text{\texttt{magnitude/}} & \text{\texttt{ogs\_local.yaml}} \\
    \text{\texttt{dataset/}} &
    \begin{gathered}
      \text{\texttt{instance.yaml}}, \\
      \text{\texttt{stead.yaml}}, \\
      \text{\texttt{scedc.yaml}}
    \end{gathered}
  \end{cases}
\end{equation}

Key operational advantages include:
\begin{itemize}
  \item \textbf{Dynamic Command-Line Composition}: Stages and
    hyperparameters are overridden dynamically at submission without
    modifying source files:
  \begin{lstlisting}[language=bash,basicstyle=\scriptsize\ttfamily]
ml_catalog_run +libpath=. picker=phasenet picker.threshold=0.2 associator=pyocto
  \end{lstlisting}
  \item \textbf{Strict Provenance Tracking}: For every execution
    run, Hydra serializes a timestamped, frozen snapshot of the
    complete resolved configuration (\texttt{.hydra/config.yaml})
    directly into the target output directory alongside generated
    catalogs. This guarantees that any published catalog entry can
    be traced back to its parameterization.
\end{itemize}

\subsection{Makefile-Driven \acs{DAG} Execution}
The macro-scale orchestration across cluster nodes is driven by a
master Makefile (\texttt{OGS/utils/Leonardo/Makefile}). The workflow
defines a \ac{DAG} governing data dependencies,
as illustrated in \cref{fig:hpc_dag}.

\begin{figure}[htbp]
  \centering
  \begin{tikzpicture}[
      >=Stealth,
      node distance=5mm and 6mm,
      every node/.style={font=\scriptsize},
      base/.style={rectangle, rounded corners=3pt, draw=black!80,
      line width=0.7pt, align=center, inner sep=4pt, minimum height=7mm},
      data/.style={base, fill=gray!10, draw=black!80, double, double
      distance=1pt, text width=26mm},
      proc/.style={base, fill=blue!8, draw=blue!80!black, text width=28mm},
      qc/.style={base, fill=teal!10, draw=teal!80!black, text width=26mm},
      arrow/.style={->, thick, draw=black!75}
    ]

    % Nodes
    \node[data] (raw) {Raw MSEED Streams\\(\texttt{WAVE\_PATH})};
    \node[proc, right=of raw] (pick) {Phase Picking\\(\ac{PhaseNet}
    / \ac{EQT})};
    \node[proc, right=of pick] (assoc) {Phase
    Association\\(\ac{GaMMA} / \ac{PyOcto})};
    \node[proc, right=of assoc] (loc) {Hypocenter
    Location\\(\ac{NonLinLoc} 1D/3D)};
    \node[proc, below=of loc] (mag) {Magnitude
    Estimation\\(\texttt{OGSLocalMagnitude})};
    \node[qc, below=of assoc] (bgma) {\acs{BPGMA} Evaluation\\(Bipartite
    Matching)};
    \node[data, below=of pick] (cat) {Final Catalog\\(Parquet / CSV)};

    % Connections
    \draw[arrow] (raw) -- (pick);
    \draw[arrow] (pick) -- (assoc);
    \draw[arrow] (assoc) -- (loc);
    \draw[arrow] (loc) -- (mag);
    \draw[arrow] (mag) -- (bgma);
    \draw[arrow] (bgma) -- (cat);

    % Layering Box
    \begin{pgfonlayer}{background}
      \node[draw=blue!40, fill=blue!2, dashed, rounded corners=5pt,
      fit=(pick) (assoc) (loc) (mag), inner sep=4pt] (pipelinebox) {};
      \node[anchor=north west,
      font=\tiny\bfseries\color{blue!70!black}] at
      (pipelinebox.north west) {Leonardo Booster Execution Pipeline};
    \end{pgfonlayer}

  \end{tikzpicture}
  \caption{\acf{DAG} of the automated \ac{HPC} seismic
    data processing pipeline. The reader should notice the
    sequential dependency
    chain where each stage consumes immutable data contracts
    produced by its predecessor,
    with \acs{BPGMA} validation integrated prior to final catalog
    export. Grounded in the
  production Makefile and Slurm orchestrators in \texttt{OGS/utils/Leonardo/}.}
  \label{fig:hpc_dag}
\end{figure}

The Makefile implements bracketed parameterized targets:
\begin{lstlisting}[language=bash,basicstyle=\scriptsize\ttfamily]
make "PhaseNet[INSTANCE,0.2]" YEAR=2024
make "PyOcto[PhaseNet,INSTANCE,0.2]" YEAR=2024
make "NLL1D[PyOcto,PhaseNet,INSTANCE,0.2]" YEAR=2024 \\
  START_DATE=20240101 END_DATE=20241231
\end{lstlisting}

This architecture prevents redundant computation: if pick
assignments for a given year already exist on Lustre, downstream
association targets execute immediately without re-triggering neural
network inference. Furthermore, symbolic link mapping (\texttt{ln
-s}) shares intermediate results across experiments, minimizing
storage consumption.

% =============================================================================
% SECTION 4: PARALLELIZATION STRATEGIES AND OPTIMIZATION
% =============================================================================
\section{Parallelization Strategies and Performance Optimization}
\label{sec:parallelization}

Achieving computational efficiency requires matching the
mathematical characteristics of each pipeline stage to the
appropriate hardware acceleration paradigm. The architecture deploys
a hybrid parallelization model spanning \ac{GPU} tensor acceleration,
shared-memory \ac{CPU} multithreading, and distributed-memory task farming.

\subsection{Hardware Resource Invariant Constraint}
To prevent \ac{CPU} core oversubscription and eliminate thread-context
switching overhead, the allocation strategy enforces the strict
conservation law:
\begin{equation}
  N_{\mathrm{tasks}} \times N_{\mathrm{cpus\_per\_task}} \equiv
  N_{\mathrm{cores}} = \num{32},
\end{equation}
where $N_{\mathrm{cores}} = \num{32}$ corresponds to the total physical
core count of the Intel Xeon 8358 processor.

\Cref{tab:thread_allocation} demonstrates how this invariant
configures varying operational execution modes:

\begin{table}[htbp]
  \centering
  \small
  \caption{Processor and thread allocation configurations under the
    core invariant constraint. The reader should notice how the
    \num{4}-task layout
    achieves $\num{1}:\num{1}$ mapping between tasks and NVIDIA Ampere A100
    \acp{GPU} with \num{8}
    \ac{CPU} cores per accelerator.
    Formulated in production launcher scripts (\texttt{LAUNCHME.sh},
  lines 41--48).}
  \label{tab:thread_allocation}
  \begin{threeparttable}
    \begin{tabularx}{\linewidth}{c c c >{\raggedright\arraybackslash}X}
      \toprule
      \textbf{Nodes} ($N_{\mathrm{nodes}}$) &
      \textbf{Threads} ($N_{\mathrm{cores}}$) &
      \textbf{\acp{GPU}} & \textbf{Target Pipeline Stage / Paradigm} \\
      \midrule
      \num{1} & \num{32} & \num{1} & Serial orchestration / Deep
      grid generation (\ac{NonLinLoc}) \\
      \num{2} & \num{16} & \num{2} & Memory-intensive association
      (large \ac{GaMMA} swarms) \\
      \num{4} & \num{8}  & \num{1} & Multi-\ac{GPU} deep learning
      inference (\ac{PhaseNet} / \ac{EQT}) \\
      \bottomrule
    \end{tabularx}
    \begin{tablenotes}[flushleft]
      \footnotesize
    \item Target \& Scope: Processor core pinning across Leonardo
      Booster compute nodes.
    \item What the reader should notice: The \num{4}-task configuration
      achieves $\num{1}:\num{1}$ mapping between tasks and NVIDIA
      Ampere A100 \acp{GPU} while
      dedicating \num{8} \ac{CPU} cores per \ac{GPU} for
      asynchronous preprocessing.
    \item Evidence Anchor: Formulated in \texttt{LAUNCHME.sh} (lines 41--48).
    \end{tablenotes}
  \end{threeparttable}
\end{table}

\subsection{Multi-GPU Neural Network Inference Optimization}
Phase picking is computationally dominant in terms of raw
floating-point operations. Accelerating inference across the \num{4}
NVIDIA Ampere A100 \acp{GPU} involves four key optimizations:
\begin{enumerate}
  \item \textbf{Asynchronous CUDA Streams and Pinned Memory}:
    Waveform arrays representing continuous three-component particle
    velocity $\mathbf{v}(t) \equiv \dot{\mathbf{u}}(t)$ are staged in
    page-locked (pinned) host memory (\texttt{pin\_memory=True}). Data
    transfers between host \ac{RAM} and \ac{GPU} HBM2e memory are executed
    via asynchronous CUDA streams (\texttt{torch.cuda.Stream}), completely
    overlapping PCIe data transfers with Tensor Core execution.
  \item \textbf{TensorFloat-32 (TF32) Execution}: Neural network
    operations utilize Ampere TF32 precision, which matches FP32
    range while delivering the throughput of half-precision matrix
    math without requiring manual loss scaling or network re-calibration.
  \item \textbf{Dynamic Window Batching}: Seismograms from all
    active network stations are assembled into contiguous 3D tensors:
    \begin{equation}
      \mathbf{B} \in \mathbb{R}^{B_{\mathrm{size}} \times 3 \times
      N_w}, \quad B_{\mathrm{size}} \in [\num{128}, \num{512}],
    \end{equation}
    saturating the \num{108} Streaming Multiprocessors (SMs) of the
    NVIDIA Ampere A100
    \ac{GPU} and sustaining $> \qty{85}{\percent}$ active compute utilization.
\end{enumerate}

\subsection{Shared-Memory Multithreading (OpenMP and Numba)}
For \ac{CPU}-bound stages---such as spatial ray tracing in \ac{NonLinLoc}, 4D
priority-queue branching in \ac{PyOcto}, and geodetic coordinate
transformations in \ac{GaMMA}---parallelism is achieved via OpenMP and
Numba multithreading:
\begin{itemize}
  \item Thread count is dynamically bounded by environment variables:
    \begin{equation}
      \text{\texttt{export OMP\_NUM\_THREADS=\$SLURM\_CPUS\_PER\_TASK}}
    \end{equation}
    \begin{equation}
      \text{\texttt{export NUMBA\_NUM\_THREADS=\$SLURM\_CPUS\_PER\_TASK}}
    \end{equation}
  \item Thread affinity is locked via OpenMP runtime directives
    (\texttt{OMP\_PROC\_BIND=close}, \texttt{OMP\_PLACES=cores}),
    binding worker cores to physically adjacent \ac{CPU} cores and
    maximizing L2/L3 cache reuse.
\end{itemize}

\subsection{Parallel I/O Optimization on Lustre}
Processing continuous seismic telemetry across years involves
millions of individual file reads. Naive I/O access patterns
generate severe metadata contention on the Lustre Metadata Targets (MDTs).

To ensure high I/O throughput:
\begin{enumerate}
  \item \textbf{Lustre File Striping}: Target directories are
    configured with striping across Object Storage Targets:
  \begin{lstlisting}[language=bash,basicstyle=\scriptsize\ttfamily]
lfs setstripe -c 4 -S 4M /leonardo_work/IscrC_AISeism/WORK/waveform
  \end{lstlisting}
    This distributes contiguous waveform blocks across \num{4} OSTs,
    enabling parallel read streams during multi-station loading.
  \item \textbf{Memory-Mapped Serialization}: Output picks and event
    associations are serialized using the Apache Parquet columnar
    format with Snappy compression, reducing disk storage footprint
    by \qty{75}{\percent} compared to text-based CSVs and accelerating
    downstream bipartite matching ingestion by an order of magnitude.
  \item \textbf{Day-Sharded SQLite Waveform Indexing}: Maintaining a
    single monolithic SQLite database (\texttt{ogsDB/.squirrel}) across
    multi-year continuous waveform archives incurs severe POSIX file
    locking latency on distributed parallel filesystems (Lustre/GPFS).
    To eliminate database write contention, the data layer partitions
    indexes into autonomous calendar-day shards
    (\path{ogsDB/YYYY/MM/DD/.squirrel}).
    Each calendar day constitutes an independent SQLite environment,
    allowing worker processes to index hundreds of days fully in parallel
    without cross-worker synchronization locks.
  \item \textbf{Node-Level Multithreading and Ephemeral Staging}: During
    full-archive indexing (\texttt{ogsdata.py}), a thread pool
    (\texttt{ThreadPoolExecutor}) distributes daily shards across all
    available physical \ac{CPU} cores ($N_{\mathrm{cores}} \coloneqq \num{32}$
    on Booster or \num{112} on \ac{DCGP}). Furthermore, when downstream phase
    pickers query waveforms, each day shard (typically $< \qty{5}{\mega\byte}$)
    is optionally staged into node-local temporary storage (\path{/tmp} or
    \texttt{tmpfs}) via \texttt{copy\_to\_tmp=True}, achieving
    \qty{100}{\percent} lock-free, low-latency
    waveform lookups during high-concurrency neural inference.
\end{enumerate}

% =============================================================================
% SECTION 5: HPC SCALABILITY ANALYSIS AND RESOURCE SYNTHESIS
% =============================================================================
\section{\acs{HPC} Scalability Analysis and Resource Synthesis}
\label{sec:scalability}

To evaluate the operational throughput of the computational pipeline
on the Leonardo supercomputer, comprehensive scaling benchmarks were
conducted across each distinct processing stage.

\subsection{Pipeline Stage Resource Allocation and Throughput}
\Cref{tab:stage_benchmarks} records the intended
resource-allocation and throughput reporting layout for a benchmark
dataset of \num{30} continuous seismic stations over one operational
calendar year (\num{365} continuous days). Its quantities require
replacement with values from versioned benchmark logs.
\begin{landscape}
  \begin{table}[htbp]
    \centering
    \small
    \caption{Unverified draft layout for resource allocation and
      throughput by pipeline stage. The reader should notice how computational
      bottlenecks shift from \ac{GPU}-dominated tensor inference in picking to
      \ac{CPU}-dominated graph and octree traversals in association. Evidence
      limitation: values represent unverified target layouts
      requiring replacement
    from archived benchmark logs.}
    \label{tab:stage_benchmarks}
    \begin{threeparttable}
      \begin{tabularx}{\linewidth}{l c c c c X}
        \toprule
        \textbf{Pipeline Stage} & \textbf{Nodes} & \textbf{Tasks} &
        \textbf{Threads} & \textbf{Accelerators} & \textbf{Empirical
        Throughput / Runtime\tnote{a}} \\
        \midrule
        0. Preprocessing \& Filtering & \num{1} & \num{4} & \num{8}
        & \ac{CPU} & $\sim
        \qty{2.5}{\minute} / \text{station-year}$ [Derived Index] \\
        1. Phase Picking (\ac{PhaseNet}) & \num{1} & \num{4} &
        \num{8} & $\num{4} \times $  NVIDIA Ampere A100 &
        $\sim \qty{42}{\second} / \text{station-year}$ ($> \num{2000} \times $
        real-time) \\
        1. Phase Picking (\ac{EQT}) & \num{1} & \num{4} & \num{8} &
        $\num{4} \times $  NVIDIA Ampere A100
        & $\sim \qty{2.1}{\minute} / \text{station-year}$ ($>
        \num{250} \times $ real-time) \\
        2. Phase Association (\ac{PyOcto}) & \num{1} & \num{1} &
        \num{32} & Multi-Core &
        $\sim \qty{15}{\second} / \text{network-day}$ ($> \num{5000} \times $
        real-time) \\
        2. Phase Association (\ac{GaMMA}) & \num{1} & \num{2} &
        \num{16} & Multi-Core & $\sim
        \qty{180}{\second} / \text{network-day}$ ($\sim \num{480} \times $
        real-time) \\
        3. Relocation (\ac{NonLinLoc}) & \num{1} & \num{1} &
        \num{32} & OpenMP & $\sim
        \qty{0.12}{\second} / \text{event}$ (Oct-Tree global search) \\
        4. Magnitude ($M_{\mathrm{L}}$) Estimation & \num{1} &
        \num{4} & \num{8} & Multi-Core &
        $\sim \qty{0.05}{\second} / \text{event}$ (Wood-Anderson simulation) \\
        5. \acs{BPGMA} Bipartite Matching & \num{1} & \num{1} &
        \num{32} & NetworkX & $\sim
        \qty{1.8}{\second} / \text{catalog-day}$ (Optimal assignment) \\
        \bottomrule
      \end{tabularx}
      \begin{tablenotes}[flushleft]
        \footnotesize
      \item[a] [Unverified Draft Value] Values are not validated
        runtime measurements in this checkout.
      \item Target \& Scope: Intended quantitative performance
        characterization of the complete computational seismic pipeline.
      \item What the reader should notice: the table separates
        workload, allocation, and throughput so a future benchmark can
        identify its computational bottleneck.
      \item Evidence requirement: versioned workload definition,
        resolved configuration, allocation record, repetitions, and
        raw timing logs.
      \end{tablenotes}
    \end{threeparttable}
  \end{table}
\end{landscape}

\subsection{Amdahl's and Gustafson's Scaling Behavior}
The scalability of the workflow across multi-node allocations
follows two complementary regimes:
\begin{enumerate}
  \item \textbf{Embarrassingly Parallel Temporal Decomposition}:
    Continuous seismic waveforms partition cleanly into calendar
    days or monthly intervals. Inter-day dependencies are limited to
    edge window overlapping ($W_O \le \qty{30}{\second}$). Distributing
    calendar days across $P$ compute nodes achieves near-linear weak scaling:
    \begin{equation}
      S(P) \coloneqq \frac{T_1}{T_P} \approx P \cdot \eta_{\mathrm{I/O}},
    \end{equation}
    where $\eta_{\mathrm{I/O}} \in [0,1]$ is the measured Lustre
    parallel-I/O efficiency for a specified workload; it must be
    estimated from benchmark logs rather than assumed.
  \item \textbf{Intra-Day Association Scaling (Amdahl's Law)}: For intensive
    seismic swarms occurring within single 24-hour windows,
    association speedup under fixed workload size is bounded by Amdahl's Law:
    \begin{equation}
      S_{\mathrm{Amdahl}}(N) \coloneqq \frac{1}{(1 - f) + \frac{f}{N}},
    \end{equation}
    where $f \in [0,1]$ is the workload-specific parallelizable
    fraction. Its value and the resulting speedup require repeated
    measurements for the chosen \ac{PyOcto} workload.
  \item \textbf{Scaled Workload Expansion (Gustafson's Law)}: When problem
    size scales proportionally with available compute resources---such as
    expanding network station density or temporal resolution while maintaining
    bounded execution time---speedup is governed by Gustafson's
    scaled speedup law:
    \begin{equation}
      S_{\mathrm{Gustafson}}(P) \coloneqq P - s(P - 1) \equiv (1 - s)P + s,
    \end{equation}
    where $s \in [0, 1]$ represents the serial runtime fraction of
    the expanded problem.
\end{enumerate}

The combination of \ac{CINECA} Leonardo's Booster architecture,
disciplined Slurm affinity binding, and Hydra-Makefile orchestration
establishes a resilient, high-throughput computational platform
capable of processing years of continuous seismic records in a matter of hours.
